\documentclass[12pt,a4paper]{article}

% --- Packages ---
\usepackage[utf8]{inputenc}
\usepackage[T1]{fontenc}
\usepackage[french]{babel}
\usepackage{geometry}
\usepackage{enumitem}
\usepackage{array}
\usepackage{fancyhdr}
\usepackage{amsmath, amssymb}

% --- Mise en page ---
\geometry{margin=2.5cm}

% --- Pied de page ---
\pagestyle{fancy}
\fancyhf{}
\fancyfoot[R]{\it{G.M.}}
\renewcommand{\headrulewidth}{0pt}
\renewcommand{\footrulewidth}{0pt}

\setlength{\parindent}{0pt}

\begin{document}

\begin{center}
    \Large \textbf{--Automatismes - QCM --} \\[5pt]
    \normalsize Aucune justification n’est demandée et une seule réponse est possible par question.
\end{center}

\vspace{0.7cm}

\textbf{Question 1 :}
\vspace{5pt}

Dans une classe de 35 élèves, 7 sont gauchers. La proportion d’élèves gauchers est :
\vspace{0.3cm}

\renewcommand{\arraystretch}{3}
\begin{tabular}{|m{3.4cm}|m{3.4cm}|m{3.4cm}|m{3.4cm}|}
\hline
\centering \vspace{-0.3cm} \textbf{A.} 0,2 & 
\centering \vspace{-0.3cm} \textbf{B.} 0,25 & 
\centering \vspace{-0.3cm} \textbf{C.} $\dfrac{1}{3}$ & 
\centering \vspace{-0.3cm} \textbf{D.} 0,4 \tabularnewline
\hline
\end{tabular}

\vspace{0.7cm}

\textbf{Question 2 :}
\vspace{5pt}

Un article coûte 160 €. Le vendeur propose une réduction de 20\%. Combien d'euros représente la réduction ?
\vspace{0.3cm}

\renewcommand{\arraystretch}{3}
\begin{tabular}{|m{3.4cm}|m{3.4cm}|m{3.4cm}|m{3.4cm}|}
\hline
\centering \vspace{-0.3cm} \textbf{A.} 28 € & 
\centering \vspace{-0.3cm} \textbf{B.} 19 € & 
\centering \vspace{-0.3cm} \textbf{C.} 32 € & 
\centering \vspace{-0.3cm} \textbf{D.} 34 € \tabularnewline
\hline
\end{tabular}

\vspace{0.7cm}

\textbf{Question 3 :}
\vspace{5pt}

Dans un collège, 15\% des élèves sont en classe de 3ème, ce qui représente 36 élèves. L’effectif total du collège est :
\vspace{0.3cm}

\renewcommand{\arraystretch}{3}
\begin{tabular}{|m{3.4cm}|m{3.4cm}|m{3.4cm}|m{3.4cm}|}
\hline
\centering \vspace{-0.3cm} \textbf{A.} 350 & 
\centering \vspace{-0.3cm} \textbf{B.} 300 & 
\centering \vspace{-0.3cm} \textbf{C.} 240 & 
\centering \vspace{-0.3cm} \textbf{D.} 280 \tabularnewline
\hline
\end{tabular}

\vspace{0.7cm}

\textbf{Question 4 :}
\vspace{5pt}

Un prix augmente de 8\%. Le coefficient multiplicateur associé à cette augmentation est :
\vspace{0.3cm}

\renewcommand{\arraystretch}{3}
\begin{tabular}{|m{3.4cm}|m{3.4cm}|m{3.4cm}|m{3.4cm}|}
\hline
\centering \vspace{-0.3cm} \textbf{A.} 0,08 & 
\centering \vspace{-0.3cm} \textbf{B.} 1,8 & 
\centering \vspace{-0.3cm} \textbf{C.} 0,92 & 
\centering \vspace{-0.3cm} \textbf{D.} 1,08 \tabularnewline
\hline
\end{tabular}

\vspace{0.7cm}

\textbf{Question 5 :}
\vspace{5pt}

Un article coûte 70 €. Après une augmentation de 10\%, son nouveau prix est :
\vspace{0.3cm}

\renewcommand{\arraystretch}{3}
\begin{tabular}{|m{3.4cm}|m{3.4cm}|m{3.4cm}|m{3.4cm}|}
\hline
\centering \vspace{-0.3cm} \textbf{A.} 70,1 € & 
\centering \vspace{-0.3cm} \textbf{B.} 77 € & 
\centering \vspace{-0.3cm} \textbf{C.} 80 € & 
\centering \vspace{-0.3cm} \textbf{D.} 81,7 € \tabularnewline
\hline
\end{tabular}

\vspace{0.7cm}

\textbf{Question 6 :}
\vspace{5pt}

Un article coûte 220 €. Après une réduction de 10\%, son nouveau prix est :
\vspace{0.3cm}

\renewcommand{\arraystretch}{3}
\begin{tabular}{|m{3.4cm}|m{3.4cm}|m{3.4cm}|m{3.4cm}|}
\hline
\centering \vspace{-0.3cm} \textbf{A.} 202 € & 
\centering \vspace{-0.3cm} \textbf{B.} 195 € & 
\centering \vspace{-0.3cm} \textbf{C.} 200 € & 
\centering \vspace{-0.3cm} \textbf{D.} 198 € \tabularnewline
\hline
\end{tabular}

\vspace{0.7cm}

\textbf{Question 7 :}
\vspace{5pt}

Après une augmentation de 40\%, un article coûte 98 €. Son prix initial était :
\vspace{0.3cm}

\renewcommand{\arraystretch}{3}
\begin{tabular}{|m{3.4cm}|m{3.4cm}|m{3.4cm}|m{3.4cm}|}
\hline
\centering \vspace{-0.3cm} \textbf{A.} 80 € & 
\centering \vspace{-0.3cm} \textbf{B.} 70 € & 
\centering \vspace{-0.3cm} \textbf{C.} 60 € & 
\centering \vspace{-0.3cm} \textbf{D.} 65 € \tabularnewline
\hline
\end{tabular}

\vspace{0.7cm}

\textbf{Question 8 :}
\vspace{5pt}

Après une réduction de 22\%, un article coûte 78 €. Son prix initial était :
\vspace{0.3cm}

\renewcommand{\arraystretch}{3}
\begin{tabular}{|m{3.4cm}|m{3.4cm}|m{3.4cm}|m{3.4cm}|}
\hline
\centering \vspace{-0.3cm} \textbf{A.} 95 € & 
\centering \vspace{-0.3cm} \textbf{B.} 100 € & 
\centering \vspace{-0.3cm} \textbf{C.} 105 € & 
\centering \vspace{-0.3cm} \textbf{D.} 98 € \tabularnewline
\hline
\end{tabular}

\vspace{0.7cm}

\textbf{Question 9 :}
\vspace{5pt}

Un prix passe de 75 € à 84 €. Le taux d’évolution, exprimé en pourcentage, est :
\vspace{0.3cm}

\renewcommand{\arraystretch}{3}
\begin{tabular}{|m{3.4cm}|m{3.4cm}|m{3.4cm}|m{3.4cm}|}
\hline
\centering \vspace{-0.3cm} \textbf{A.} 10\% & 
\centering \vspace{-0.3cm} \textbf{B.} 16\% & 
\centering \vspace{-0.3cm} \textbf{C.} 12\% & 
\centering \vspace{-0.3cm} \textbf{D.} 14\% \tabularnewline
\hline
\end{tabular}

\vspace{0.7cm}

\textbf{Question 10 :}
\vspace{5pt}

Un prix subit deux augmentations successives de 6\% puis de 10\%. Le taux d’évolution global, exprimé en pourcentage, est :
\vspace{0.3cm}

\renewcommand{\arraystretch}{3}
\begin{tabular}{|m{3.4cm}|m{3.4cm}|m{3.4cm}|m{3.4cm}|}
\hline
\centering \vspace{-0.3cm} \textbf{A.} 16,6\% & 
\centering \vspace{-0.3cm} \textbf{B.} 15,6\% & 
\centering \vspace{-0.3cm} \textbf{C.} 17,2\% & 
\centering \vspace{-0.3cm} \textbf{D.} 16\% \tabularnewline
\hline
\end{tabular}

\end{document}
